\chapter{Uncertain relationships, transported policies and the information to audit}
\label{exec17:s7}

\section{A banking relationship is an input to a policy's meaning}
\label{exec17:s7-question}

Consider an outreach proposal intended to improve repayment through a customer's
business contacts. A contact may share information, coordinate a payment or pass
on an offer. The bank has candidate relationships assembled from several files,
and a pilot population differs from the proposed rollout population. The question
is the incremental value of changing the outreach allocation. It requires a
definition of the relationships through which the action can operate, the joint
allocation of actions, the population receiving the benefit and the observations
that establish those quantities. An accurate estimate conditional on an invented
network can answer the wrong banking question.

There are three distinct objects. An identity link says that two records refer
to the same entity. A relationship link says that two distinct entities have a
specified connection. An influence link asserts that an action can affect an
outcome through that connection. A common telephone number can suggest the first
or second object without establishing the third. Conversely, a one-to-one
customer match does not establish a one-to-one guarantee network: a firm can have
many guarantors. The observation methods in Chapter~\ref{r3:observation} and the
allocation methods in Chapter~\ref{v13:sw-networks} address parts of this problem.
Here the issue is their composition: uncertainty in the first objects changes
the exposure, target population and monetary value used by the later methods.

Let $G$ denote the scientifically relevant relationship graph, $L$ an identity
assignment, $R$ the observation indicators and $Z$ a pretreatment population
stratum. For an action vector $\mathbf a$, an exposure mapping
$h_i(G,\mathbf a)$ specifies which part of that vector can affect customer $i$.
A policy $g$ is a joint probability law for feasible action vectors, conditional
on information available at its decision time. It need not factor into
independent customer-level coins. A useful population value is
\begin{equation}
 V_*(g;M)=\E_{M,*}\!\left[
   \sum_{\mathbf a}g(\mathbf a\mid\mathcal F_{-})
     \{Y_i(h_i(G,\mathbf a))-C_i(\mathbf a)\}\right],
 \label{exec17:s7-target}
\end{equation}
where $M$ specifies the outcome, graph and observation law, $*$ denotes the
target population, and costs are allocated once to the stated beneficiary or
population denominator. For a finite population the outer expectation includes
a uniformly selected target member. Other population weights require an
explicit replacement. A policy contrast subtracts the old value and the
incremental investment cost.

This notation does not identify the value. It reveals the inputs that must
remain common when comparing software: the observation law for $L,G,R$;
the exposure mapping; the full action law; the target population; and the
currency and horizon. If several admissible models agree on the supplied
observations, the scientifically warranted result can be a set of values.
The examples below are project derivations with exact, hypothetical monetary
inputs. They establish semantic and information limits before a fitted
estimator or an expensive bank experiment is chosen.~\citep{exec17_s7_math}

\section{Separate edge probabilities do not determine policy exposure}
\label{exec17:s7-overlap}

Start with disjoint candidate neighborhoods. Each has one focal beneficiary and
two distinct actionable contacts. Let $U,V\in\{0,1\}$ indicate whether the
respective influence relationships are real. Assume their marginal probabilities
are both $1/2$. The statement concerns a population distribution of actual
neighborhoods; it is not a claim that a particular customer's graph changes
randomly at every decision. Conditional on a fixed neighborhood, independently
assign each contact an action with probability $p$, independently of $(U,V)$.
The focal beneficiary benefits if at least one real contact receives an action.

Inclusion and exclusion of the two successful-contact events gives
\begin{equation}
 q(p\mid U,V)=p(U+V)-p^2UV.
 \label{exec17:s7-exposure}
\end{equation}
The first term counts both possible contacts; the second removes double counting
when both relationships are real and both contacts are treated. Write
$\theta=\Pr(U=V=1)$. The supplied margins imply
\begin{equation}
 \begin{split}
  &\Pr(1,1)=\Pr(0,0)=\theta,\\
  &\Pr(1,0)=\Pr(0,1)=\tfrac12-\theta,
  \qquad 0\leq\theta\leq\tfrac12,\\
  &q(p)=p-p^2\theta.
 \end{split}
 \label{exec17:s7-sharp-family}
\end{equation}
Every member of this family is a coherent joint graph law. At $\theta=0$ every
beneficiary has exactly one real contact; at $\theta=1/2$ half have two and half
have none. Both populations have the same separate edge probabilities. They
produce different saturation of outreach because the second contact has less
incremental value when the first already reaches the beneficiary.

Suppose a successful exposure contributes USD 100. An action costs USD 20
for each of the two potential contacts, including contacts whose influence
relationship is absent. Raising $p$ from $1/5$ to $3/5$ increases expected contact
expense by $2(20)(2/5)=16$ dollars per focal beneficiary. Therefore
\begin{equation}
 \Delta q=\tfrac25-\tfrac8{25}\theta,
 \qquad \Delta V_{\mathrm{before\ investment}}=24-32\theta.
 \label{exec17:s7-graph-value}
\end{equation}
An additional USD 16 investment per beneficiary gives
\begin{equation}
 v(\theta)=8-32\theta,\qquad
 \{v(\theta):0\leq\theta\leq\tfrac12\}=[-8,8].
 \label{exec17:s7-graph-set}
\end{equation}
The bounds are sharp because the endpoint graph distributions in
\eqref{exec17:s7-sharp-family} attain them. Enumerating the four graph states
and four action states reproduces the expectation. These are exact structural
calculations; the interval is not a confidence interval from a sampled bank.
\citep{exec17_s7_math}

Completing the two edges independently selects $\theta=1/4$ and reports zero
net value after investment. Independence of graph edges is additional scientific
information, distinct from independent allocation of actions. The latter was
specified as a policy; the former was not supplied by the edge margins. A
calibrated probability of each edge does not alone calibrate their joint
occurrence. This matters even when the outcome model and monetary inputs are
known exactly.

It also changes what should be audited. Suppose indefinitely many independent
neighborhoods are audited, but each audit randomly chooses only one of the two
edges and never records both. Under every $\theta$ in the family, the selected
edge is a fair Bernoulli observation. With this observation design the complete
audit-data law is identical across all $\theta$. More separate-edge labels
cannot distinguish the endpoint policy values. Joint observations, a justified
dependence restriction or a different policy are needed. If outcomes from a
rich experiment were also supplied, they could contain more information; that
is a different experiment from this deliberately limited audit design.

\section{The joint action law is a substantive alternative}
\label{exec17:s7-coordination}

Uncertain relationships need not force the bank to choose a graph completion.
Sometimes it can change the action design to reduce reliance on the unknown
graph moment. Let the two contact actions have equal marginal probability $p$
and joint probability $r=\Pr(A_1=A_2=1)$. Suppose the pair of actions is independent
of the graph. Its feasible distribution is
\begin{equation}
 \begin{split}
  &\Pr(1,1)=r,\quad \Pr(1,0)=\Pr(0,1)=p-r,\\
  &\Pr(0,0)=1-2p+r,\quad
      \max(0,2p-1)\leq r\leq p,\\
  &q(p,r)=p-r\theta.
 \end{split}
 \label{exec17:s7-action-law}
\end{equation}
The same separate action probabilities can thus represent independent coins,
all-or-none outreach or a plan that disperses actions across contact pairs.
Their expected action expense is the same, but their policy exposure differs.

For a matched comparison, fix the old action law at $p_0=r_0=1/5$: both
contacts receive outreach with probability $1/5$, otherwise neither does. All
new plans have $p_1=3/5$. Retain the USD 100 exposure benefit, USD 20 contact
expense and USD 16 incremental investment. The value becomes
\begin{equation}
 v(r_1;\theta)=8-100(r_1-\tfrac15)\theta.
 \label{exec17:s7-coordinated-value}
\end{equation}
Table~\ref{exec17:s7-policy-table} compares three feasible new plans against
this same old law. Its baseline differs from the independently randomized old
law in \eqref{exec17:s7-graph-set}; keeping the reference explicit prevents
unmatched claims of improvement.

\begin{table}[htbp]
\centering\small
\begin{tabular}{p{0.26\textwidth}p{0.15\textwidth}p{0.23\textwidth}p{0.18\textwidth}}
\toprule
New allocation & $r_1$ & Net value & Sharp value set\\
\midrule
Independent contact coins & $9/25$ & $8-16\theta$ & $[0,8]$\\
Both contacts together & $3/5$ & $8-40\theta$ & $[-12,8]$\\
Dispersed contact pairs & $1/5$ & $8$ & $\{8\}$\\
\bottomrule
\end{tabular}
\caption{Exact hypothetical USD values per focal beneficiary, with the same
old allocation and the same marginal action expense.}
\label{exec17:s7-policy-table}
\end{table}

The dispersed new plan assigns both actions with probability $1/5$ and each
single-action configuration with probability $2/5$; it never assigns neither.
Its overlap probability equals the old one, so the unknown $\theta$ cancels
from the contrast. The plan has a positive USD 8 value throughout the supplied
graph class. The conclusion relies on the stipulated saturation response and
linear contact expense. An action interaction in cost or a benefit depending
on which contact is reached would change it.~\citep{exec17_s7_math}

Feasibility is equally substantive. Disjoint actionable pairs allow these
pair distributions to be implemented jointly. In overlapping neighborhoods,
one contact can belong to several pairs, and separately specified pair laws
may not extend to a common global action law. The product must then solve or
verify the global allocation problem before presenting the policy value.
It must also preserve a common action clock: a sequential rule that reacts
to the first contact's response is a different intervention.

This example offers a route beyond requesting more labels. Information
acquisition and intervention design are alternatives to compare. Where
coordination is feasible, changing the plan can remove the particular unknown
moment at unchanged expected contact expense. Where coordination is infeasible,
the original value set and the need for joint relationship information remain.
A software field containing only a scalar allocation probability would lose
this distinction unless its meaning explicitly fixes the entire action law.

\section{A population can change through its graph}
\label{exec17:s7-topology}

Transporting a pilot result requires invariant causal mechanisms or a justified
formula combining the available populations. Selection diagrams express such
invariance claims; the pretreatment admissibility result of
\citet[Theorem 2 and Corollary 1]{pearltransport2014} allows target averaging of
an invariant conditional intervention response under its graphical conditions.
Equal observed customer characteristics do not establish that the relevant
conditional response is invariant when relationship structure changes.

An exact finite example makes the distinction visible. Take eight customers
with the same recorded covariates. At the source they form four disjoint
undirected pairs, so everyone has degree one. At the target, four customers
are isolated and four form a cycle, so half have degree zero and half degree
two. Both graphs have four edges and mean degree one. Independently assign
each customer an action with probability $p$. An individual's stipulated
benefit is USD 100 if at least one of its neighbors receives an action; its
own action costs USD 20. There is no direct own-action benefit, and no
influence beyond immediate neighbors.

For degree $d$, the exposure probability is $1-(1-p)^d$. Averaging over these
two actual graphs gives
\begin{equation}
 q_s(p)=p,\qquad q_*(p)=\tfrac12\{1-(1-p)^2\}
                    =p-\tfrac12p^2.
 \label{exec17:s7-topology-means}
\end{equation}
When $p$ rises from $1/5$ to $3/5$, the source exposure gain is $2/5$,
whereas the target gain is $6/25$. The expected own-action expense increases
by USD 8 in both populations. Net benefits before a USD 24 investment are
therefore USD 32 and USD 16; after investment the decisions have values
$+8$ and $-8$ respectively. Global enumeration of all $2^8$ action vectors
on each graph verifies these population means.~\citep{exec17_s7_math}

The response mechanism was held fixed; the graph distribution changed.
Balancing customer covariates or matching mean degree cannot reproduce the
nonlinear saturation in \eqref{exec17:s7-topology-means}. Degree distribution
would suffice for this particular supplied outcome law, but the source has
no degree-zero or degree-two observations. Adding degree to a transport
regression therefore reveals a support failure; it does not establish the
missing response surface. The exact target calculation used an explicit
structural saturation law beyond nonparametric source evidence.

For more general responses even degree is inadequate. Direction, weighted
connections, shared neighbors or actions outside the measured graph may affect
exposure. Which graph summaries can be transported depends on a sufficient
exposure mapping and the invariances of the outcome law. A claim of unchanged
graph semantics should specify the relationship type, observation coverage and
availability clock. A cycle observed after outreach can itself be an effect of
outreach and cannot silently replace the pretreatment influence graph.

This establishes a concrete comparison requirement. A transport package that
accepts a covariate table may still be useful after a valid graph reduction.
Its table interface does not justify that reduction. A network package that
accepts a target graph may compute the stipulated target exposure accurately,
but it does not validate the invariant response mechanism or recover omitted
relationships. The uncertainty and support diagnostics must accompany both
operations.

\section{A pooled graph audit may be inadequate for target rollout}
\label{exec17:s7-pooled}

Return to disjoint two-contact neighborhoods and independent action coins at
$p_0=1/5$ and $p_1=3/5$. Suppose an observed pretreatment stratum $Z\in\{0,1\}$
has source weights $(3/4,1/4)$ and target weights $(1/4,3/4)$. Within each
stratum the two edge margins are $1/2$. Write their overlap probabilities
as $\theta_0,\theta_1\in[0,1/2]$. Assume these conditional joint graph laws
are invariant between populations. Benefits per successful exposure are
stipulated at USD 100 and USD 200 respectively. Incremental contact expense
is USD 16 and investment USD 32 per focal beneficiary in either stratum.
Consequently
\begin{equation}
 v_0=-8-32\theta_0,\qquad v_1=32-64\theta_1,
 \qquad v_*=22-8\theta_0-48\theta_1.
 \label{exec17:s7-stratum-values}
\end{equation}
The target coefficients combine stratum-specific exposure sensitivity,
benefit and population weight. They show why the second stratum is much more
important to the rollout decision than its source prevalence suggests.

Suppose the available source audit reports the exact pooled joint moment
\begin{equation}
 \tfrac34\theta_0+\tfrac14\theta_1=\tfrac14.
 \label{exec17:s7-pooled-constraint}
\end{equation}
It does not release joint labels by stratum. This is an information restriction:
the analyst has the stratum proportions, separate edge margins and the pooled
overlap, together with the stipulated response law. Unrestricted source
experiments with stratum-specific outcomes are not additionally supplied.
They could identify information absent from this example.

Combining the pooled constraint with the probability bounds yields
\begin{equation}
 \theta_1=1-3\theta_0,\qquad
 \tfrac16\leq\theta_0\leq\tfrac13,
 \qquad v_*=136\theta_0-26.
 \label{exec17:s7-pooled-reduction}
\end{equation}
Thus the sharp target value set is
\begin{equation}
 \mathcal V_*=[-\tfrac{10}{3},\tfrac{58}{3}].
 \label{exec17:s7-pooled-set}
\end{equation}
For the lower endpoint take $(\theta_0,\theta_1)=(1/6,1/2)$; for the upper
take $(1/3,0)$. Each stratum then has the coherent four-state graph law in
\eqref{exec17:s7-sharp-family}, and both populations obey the same supplied
pooled audit. Rational mixtures of disjoint neighborhoods realize these laws.
The bounds are attained, rather than obtained by unconstrained independent
choices of impossible stratum parameters.~\citep{exec17_s7_math}

Assigning $\theta_0=\theta_1=1/4$ reports USD 8. It satisfies the pooled
constraint but introduces homogeneous overlap without evidence. More precise
measurement of the same pooled moment does not distinguish the endpoint
worlds. This is different from the topology example: conditional graph
invariance was granted here, yet the source information was aggregated in a
way that loses what target weighting needs.

The practical remedy is not to discard the pilot. Preserve the pooled
constraint and collect or release the missing stratified joint information.
If privacy or custody limits prevent release of individual joint labels, an
approved stratum-level sufficient audit statistic may resolve this particular
moment. If a stratum is unobserved, defend a structural restriction or retain
the value set. Replacing the target weights with the source weights would
answer a different population question.

There is a related software hazard. A workflow can report a correct source
graph posterior and a correct target reweighting routine while using them
incompatibly. If the posterior has already collapsed the stratum-specific
parameters to one pooled quantity, later reweighting cannot recreate their
joint feasible set. The proposed representation should retain the constraints
on $(\theta_0,\theta_1)$ and apply the target value to that set. This is a
small linear optimization in the example, but the essential requirement is
scientific preservation of the joint information.

\section{Completion uncertainty and sampling uncertainty answer different questions}
\label{exec17:s7-uncertainty}

The intervals above express compatibility with supplied exact information.
They are not repeated-sampling error bars. When the graph and observation
model are estimated, at least three uncertainty sources must be distinguished:
which completions remain admissible, uncertainty conditional on a completion,
and numerical error in averaging action assignments. More Monte Carlo draws
reduce the last component without resolving the first two.

For an illustrative model-based analysis, suppose the two endpoint graph
completions in \eqref{exec17:s7-graph-set} have posterior weights $1/2$ each.
Conditional value means are $-8$ and $8$, and conditional variances are each
one square dollar. The law of total variance gives
\begin{equation}
 \Var(v\mid O)=
 \E\{\Var(v\mid G,O)\mid O\}
 +\Var\{\E(v\mid G,O)\mid O\}=1+64=65.
 \label{exec17:s7-total-variance}
\end{equation}
A four-point construction with independent conditional noise $\{-1,1\}$
checks the arithmetic. Reporting the variance one after selecting a graph
omits the between-completion variation. Computing a value at the mean graph
also need not equal averaging the value: \eqref{exec17:s7-exposure} depends
on the product $UV$.~\citep{exec17_s7_math}

These posterior weights are supplied model statements. They do not establish
frequentist coverage, a justified prior or authentic probabilities for bank
relationships. The joint Bayesian linkage construction of \citet{guha2022}
illustrates how outcome information can inform uncertain matching in a common
likelihood. Its no-duplicate one-to-one linkage, no-interference SUTVA and
ignorability assumptions must remain attached. It does not supply a theorem
for replacing $G$ by an arbitrary banking influence network. Adding a network
outcome model changes the likelihood and the causal assumptions and needs its
own derivation and evaluation.

A different route is a sampling confidence region for the graph moments.
Suppose a procedure produces $\mathcal C_n$ with
$\Pr\{(\theta_0,\theta_1)\in\mathcal C_n\}\geq1-\alpha$ under a declared
truthful audit design. Mapping this region through
\eqref{exec17:s7-stratum-values} gives
\begin{equation}
 \left[\inf_{\boldsymbol\theta\in\mathcal C_n}v_*(\boldsymbol\theta),
       \sup_{\boldsymbol\theta\in\mathcal C_n}v_*(\boldsymbol\theta)\right].
 \label{exec17:s7-region-image}
\end{equation}
It covers the true value whenever the region contains the true moments,
provided the structural response, target weights and monetary inputs are
indeed known as stipulated. If these are estimated too, a joint argument
must include their error. Separate conditional intervals on arbitrarily
selected graphs do not prove graph retention.

For simultaneous stratum intervals $[L_z,U_z]$, intersect the rectangle with
any warranted pooled constraint and the probability bounds. Without the
pooled restriction, the image is
\begin{equation}
 [22-8U_0-48U_1,\;22-8L_0-48L_1].
 \label{exec17:s7-rectangle}
\end{equation}
With the restriction, optimize over the intersection; independently selecting
both pessimistic endpoints may be infeasible. An empty intersection is a
diagnostic of incompatible inputs or sampling variation under a misused exact
constraint, not permission to drop an inconvenient audit. A noisy pooled
estimate requires its own uncertainty rather than an exact equality.

Sensitivity regions, posterior mixtures and confidence regions can all be
useful. Their interpretations differ. The product should state whether a
negative lower endpoint reflects unresolved graph information, sample
precision, uncertain transport or a numerical enclosure. Each has a different
remedy. Reporting one generic uncertainty number would hide the next decision.

\section{Audit the joint unit that the value requires}
\label{exec17:s7-audit}

Return to $v(\theta)=8-32\theta$. A truthful whole-neighborhood audit records
both relationship indicators and hence $B=UV$. Under independent draws of
eligible neighborhoods, $B_1,\ldots,B_n$ are independent Bernoulli variables
with mean $\theta$. This is a new observation design. It is not justified by
calling the two edge labels independent when they concern the same customer.
Shared contacts, common ownership errors or calendar-specific adjudication
can also violate independence between nominal neighborhoods.

For completeness, an elementary concentration argument yields a prospective
design bound. If $B\in\{0,1\}$, the second derivative of its centered log
moment-generating function is the variance of a tilted Bernoulli variable,
at most $1/4$. The function and its first derivative vanish at zero, so
integrating this bound twice gives
\begin{equation}
 \log\E\exp\{t(B-\theta)\}\leq t^2/8.
 \label{exec17:s7-mgf}
\end{equation}
Independence, an exponential Markov bound and minimizing over $t>0$ give
$\Pr(\widehat\theta-\theta\geq r)\leq\exp(-2nr^2)$; the lower tail has
the same bound. Consequently a two-sided radius is
\begin{equation}
 r_n=\sqrt{\frac{\log(2/\alpha)}{2n}},\qquad
 \Pr\{|\widehat\theta-\theta|\leq r_n\}\geq1-\alpha.
 \label{exec17:s7-radius}
\end{equation}
Intersect the interval with $[0,1/2]$. A lower value certificate is
\begin{equation}
 \underline v_n=8-32\min(\tfrac12,\widehat\theta+r_n).
 \label{exec17:s7-certificate}
\end{equation}
On the concentration event this bound is no greater than the true value.
The construction depends on the specified truthful neighborhood law; it does
not validate adjudication, inclusion probabilities or network independence.

Suppose design information separately supports $\theta\leq1/8$. On that same
event, $\widehat\theta+r_n\leq\theta+2r_n$, so strict positive-value
certification is assured if $2r_n<1/8$. At $\alpha=0.05$ this is
\begin{equation}
 n>128\log40\simeq472.1766.
 \label{exec17:s7-sample-size}
\end{equation}
Thus 473 independent truthful neighborhood audits are sufficient under these
conditions. The radius is approximately $0.062446$; direct high-precision
checks distinguish 472 from 473. This is a conservative prospective statement:
conditional on the stipulated parameter restriction and sampling law, the
concentration event has probability at least $0.95$ and implies a positive
certificate. No audits have been sampled for this calculation. It neither
promises that an arbitrary observed dataset passes nor estimates audit cost.
\citep{exec17_s7_math}

If only one edge can ever be audited per neighborhood, $B$ is never observed.
The earlier nonidentification remains even with infinitely many such labels.
If whole-neighborhood inclusion is unequal but known, a design-based weighted
joint-moment estimate may be possible; its boundedness and variance must use
that inclusion law. If some relevant neighborhoods have zero joint inclusion,
an extrapolation restriction or a partial-identification result is necessary.
Untruthful joint labels require an error model or independent gold information;
larger $n$ alone does not correct systematic error.

This gives a practical audit question for the banking owner: can the audit
establish the two scientifically defined influence links together, for the
population and information time that the policy uses? Auditing two customer
names or confirming legal ownership may be easier and still answer a different
question. The costed experiment should be designed around the missing moment,
not around the most convenient field to label.

\section{Allocate audit effort to the transported decision}
\label{exec17:s7-allocation}

The stratified example gives a further choice before collecting expensive
labels. Suppose truthful independent whole-neighborhood audits are available
within each stratum, with fixed sizes $n_0,n_1>0$. Estimate the two overlaps
by sample means and use
\begin{equation}
 \widehat v_*=22-8\widehat\theta_0-48\widehat\theta_1,
 \quad
 \Var(\widehat v_*)=
 \frac{64\theta_0(1-\theta_0)}{n_0}
 +\frac{2304\theta_1(1-\theta_1)}{n_1}
 \leq\frac{16}{n_0}+\frac{576}{n_1}.
 \label{exec17:s7-audit-variance}
\end{equation}
The variance formula assumes independent strata and independent neighborhood
draws within each. It treats the other components of value as known, rather
than silently omitting their statistical error. The last inequality is a
worst-case Bernoulli variance bound; it is a design criterion, not an observed
standard error or a coverage proof.

For a fixed number $n=n_0+n_1$ of equally priced audits, a square expansion or
Cauchy--Schwarz gives
\begin{equation}
 \left(\frac{16}{n_0}+\frac{576}{n_1}\right)(n_0+n_1)
 \geq(4+24)^2=784.
 \label{exec17:s7-audit-optimum}
\end{equation}
Equality requires $n_1/n_0=6$. At 700 audits the allocation $(100,600)$ has
bound $28/25=1.12$ square dollars. Source-proportional allocation $(525,175)$
has bound $1744/525\simeq3.322$. Both buy 700 joint neighborhood labels. The
gain is due to the target value's coefficients, not a claim that one method
has better actual bank coverage.~\citep{exec17_s7_math}

Prices can change the design. Let a stratum-zero audit cost one budget unit
and a stratum-one audit four. With budget $C=n_0+4n_1$,
\begin{equation}
 \left(\frac{16}{n_0}+\frac{576}{n_1}\right)(n_0+4n_1)
 \geq(4+48)^2=2704,
 \qquad \frac{n_1}{n_0}=3\ \text{at equality}.
 \label{exec17:s7-priced-optimum}
\end{equation}
At $C=1820$, integer allocation $(140,420)$ attains bound $52/35$, whereas
source-proportional $(780,260)$ has bound $436/195$. The expensive stratum
still merits more audits because its coefficient in transported value is six
times larger in magnitude. Budget units here are hypothetical relative prices,
not measured bank invoices.

This calculation deliberately uses a simple variance criterion. A simultaneous
concentration interval could instead allocate tail probabilities and sample
sizes to control $8r_{n_0}+48r_{n_1}$. A decision rule with asymmetric losses
could optimize expected regret under an explicit probability model. The exact
pooled constraint may also permit a more efficient estimator than the direct
two-mean estimator above. Each is a different design problem; the bound in
\eqref{exec17:s7-audit-optimum} is optimal only for its stated criterion and
information use. Fixed design followed by fresh audits keeps this choice
separate from favorable sampling outcomes.

The substantive conclusion is that audit size should follow the estimand and
the target composition. A large pooled source audit can leave a decision
unresolved, while a smaller targeted joint audit may supply the missing
information. Whether that audit is worth buying still needs actual prices,
feasible coordination, defended outcome values and a bank loss threshold.

\section{Software composition must preserve the observation assumptions}
\label{exec17:s7-software}

The existing catalog identifies software that can supply useful pieces.
The proposed integration decision is narrower: admit each piece for a
declared subproblem and preserve the unresolved quantities at the boundary.
The saved source versions through the 7 October 2026 cutoff support the
following comparisons. They are technical-source findings, not executed
package trials.~\citep{exec17_s7_paths}

\subsection{Linkage output should retain the joint information needed downstream}

The inspected Splink README describes a Fellegi--Sunter-based probabilistic
linkage workflow, pair predictions and clustering into estimated identities.
It says multiple fields that are not highly correlated are favorable inputs.
Its saved documentation snapshot is content-hashed; it is not assumed to
belong to a subsequently retrieved mutable branch HEAD. Performance statements
in that README remain project-reported claims.~\citep{sw_splink}

An original Python adapter is a candidate for identity resolution. For the
present policy it must preserve candidate pairs, comparison information,
blocking coverage, estimated probabilities and the clock of the fields.
A thresholded cluster can be a proposed identity assignment, but should not
be relabeled as a verified influence graph. Pair probabilities alone cannot
recover the overlap moment in \eqref{exec17:s7-sharp-family}. A coherent joint
linkage model or sensitivity class must also respect constraints such as
one-to-one identity matching where scientifically appropriate. Independent
pair completion can violate those constraints.

Guha, Reiter and Mercatanti's joint outcome/linkage likelihood offers a
substantive alternative to fixing matches before inference. Their saved
publisher paper is specifically about duplicate-free bipartite files under
one-to-one matching and causal no-interference assumptions. Outcome feedback
can alter linkage uncertainty within that model.~\citep{guha2022} For a network
policy, such feedback is a research starting point: the exposure depends
jointly on several links and the outcome law permits interference. A direct
port of the one-to-one causal likelihood would impose the wrong causal model.
Preserving the original linkage engine and developing the new joint model
explicitly is preferable to claiming that a language translation supplies the
missing argument.

\subsection{Missingness recovery and graph discovery require separate admission}

The missing-data discovery study of \citet{tu2019} assumes no confounding or
selection bias relative to its substantive variable set, causal Markov and
faithfulness conditions, missingness indicators that are not causes of
substantive variables, faithful observability, no specified deterministic
relations between missingness indicators and no self-masking missingness.
These restrictions matter if unpaid balances are selectively unrecorded or
account closure changes later actions. A value being absent from a table
does not establish the source's missingness graph assumptions.

In the pinned causal-learn 0.1.4.8 source, the PC entry point dispatches
\texttt{mvpc=True} to its missing-value path; a supplied Fisher-Z test is
replaced by the missing-value variant and the correction name is passed
downstream. The observed commit is \texttt{0dacacf36390} from 4 September
2026.~\citep{v12_sw_causallearn} This establishes dispatch semantics. It
does not establish a theorem for arbitrary nonrandom missingness, unknown
identity links or dependent customer networks. The proposed Python bridge
must require a source-compatible missingness and testing regime, retain the
correction choice and preserve graph ambiguity. The downstream test and
correction paths need their own inspection and runtime comparison.

Agarwal and Singh's corruption-aware route uses a different information
structure: bounded approximately low-rank signal, mean-zero subexponential
measurement errors independent across rows, and a specified masking model
that permits within-row dependence while requiring independent rows.
\citep[Section 5.1, Assumptions 5.1--5.3]{corrupt2022} Repeated measurements
can supply useful structure for selected corrupted covariates. They do not
establish truthful missing edges or independent rows when several table rows
share a customer, household or guarantee. A native reconstruction of the
algorithm would retain the same scientific restrictions. Before adopting
this route, define what constitutes a row and show how the bank's observation
law supplies the required repetition and independence.

These are substantive alternatives to blind imputation. A compatible
missingness model can identify a target; a low-rank structure can supply
additional information; a sensitivity class can preserve unresolved values;
and an audit can observe the missing joint unit. Their outputs should be
compared on a common target. Completing a table first and selecting the
best-looking causal output afterwards would hide which information supplied
identification.

\subsection{A sensitivity parameter should correspond to the unresolved mechanism}

The pinned sensemakr 0.1.6 documentation describes regression sensitivity in
terms of partial $R^2$ quantities and observed-covariate benchmarks. Its
\texttt{kd} and \texttt{ky} arguments specify how strongly a hypothetical
unobserved confounder relates to treatment and outcome compared with a chosen
benchmark. They are assumptions supplied to the analysis. The manual also
distinguishes coefficient, standard-error and degrees-of-freedom inputs.
\citep{v12_sw_sensemakr}

An original R adapter is a useful proposed report for a supported regression
question. It cannot convert a benchmark choice into evidence about the
missing overlap $\theta$ or about an invariant transport mechanism. In the
present example the sensitivity parameter has a direct joint-event meaning,
an explicit feasible set and a monetary slope. Mapping it to a partial
$R^2$ restriction would require an outcome/assignment model and a derivation;
the package name supplies neither. Conversely, even knowing every edge may
leave hidden-confounding sensitivity relevant for an observational outcome
model. The two reports can coexist with their separate targets and inputs.

\subsection{Network and group engines should receive a complete causal request}

The inspected MossSpider \texttt{NetworkTMLE} source at the July 2022 commit
documents direct-contact stochastic allocation, Monte Carlo averaging and
complete node attributes, with no missing-data procedure. Its unweighted
adjacency, degree-restriction background population and directed-influence
convention need explicit semantic treatment.~\citep{v13_sw_mossspider}
The earlier network chapter develops those individual findings. Their new
consequence here is a composition boundary: a completed graph can be an
input to a conditional evaluation, but it cannot silently close the outer
graph-value set or qualify the uncertainty after completion.

A proposed original Python wrapper should first compare conditional values
on every small graph completion against exact action enumeration. A separate
outer calculation can then map a scientifically warranted joint graph region
through those conditional values. If stochastic integration is used, its
numerical error must be bounded or recorded for the extremal completions;
a spurious numerical minimum is not a scientific bound. The coordinated
pair policy also needs support for its actual joint allocation. A routine
documented only for independent allocation cannot accept the same scalar
$p$ and claim to have evaluated that coordinated plan.

Inferference 1.0.3 supplies grouped formulas, allocation probabilities and an
IPW path with several propensity-model choices in the saved May 2025 source.
The default logit integrand permits one random effect, the code requires at
least two groups, and the prose/signature disagreement about
\texttt{integrate\_allocation} requires an explicit recorded argument.
\citep{v13_sw_inferference} None of these checks establishes independent
neighborhoods or a feasible global joint action law. A grouped R adapter is
a candidate when disjoint groups, the target allocation and the sampling law
really coincide. Crossing guarantees or common audit errors require another
dependence argument. The complete integrands, variance calculations and
numerical tolerances remain to be qualified.

Original-language adapters retain mature upstream implementations and make
these boundaries easier to inspect. A shared request must retain the graph
or graph class, identity map, exposure definition, target weights, full
allocation law, observation assumptions and distinct uncertainty components.
The observed MIT text for MossSpider, GPL declaration for inferference and
GPL-3 declaration for sensemakr inform a rights review; they do not clear
every dependency or distribution.~\citep{exec17_s7_paths} A native port is
deferred until semantic and numerical parity plus the complete rights basis
are established. No evidence here supports a universal missing-data,
transport and interference engine assembled merely by chaining package calls.

\section{A public and generated comparison before costly bank experiments}
\label{exec17:s7-prototype}

A narrow prototype can answer useful questions with public structure and
explicit generated truth. A permitted public graph or relational table can
provide heterogeneous degree, group sizes and covariate imbalance. Its
relationships should retain their published meaning: a public transaction
edge is not thereby a verified repayment influence link. Generated graph
truth, action laws, response mechanisms, linkage corruption and missingness
can then be added transparently. This supplies a controlled test of the
proposed composition while avoiding a claim that public records reveal the
bank's counterfactual outcomes.~\citep{exec17_s7_design}

The first comparison should use the exact finite examples in this chapter.
Admit the identity and exposure translation, matched policy baseline,
population denominator and cost accounting before testing any statistical
performance. Independent-edge and coherent-joint completions should agree
only where the supplied graph law warrants it. Independent and coordinated
allocations should reproduce their distinct values. Target reweighting must
retain the pooled constraint and reveal support failure in the topology
example. Negative controls should include one-edge-only audits, wrong
relationship direction, shared-neighborhood adjudication errors and an
unsupported missingness regime.

The next stage can register fresh generated mechanisms and audit samples,
fixed versions, all scenarios, repetitions, numerical tolerances and stopping
rules. Compare hard matching, a justified joint probabilistic model, partial
identification and targeted whole-neighborhood audits. Compare original
engines only where they evaluate the same target and action law. Report
graph-region retention, target error, value-set containment, finite
availability, interval width, decision loss and numerical integration error
separately. Include failures and refusals in the request denominator. The
unit for uncertainty must follow the generation and audit law; repeating
several requests on one shared graph does not create independent worlds.

The software decision then has concrete evidence requirements. Adopt an
original engine for a supported subproblem after semantic, rights and
runtime checks. Integrate it through a preserving adapter when the common
request can express its restrictions. Develop a new joint model when the
scientific composition is absent. Retain a bounded answer or refuse a
specific unsupported request when information is insufficient. Porting is
an engineering option after parity, not a source of identification.

Actual bank experiments remain a later investment decision. They need an
owned influence definition, truthful joint audits, a feasible action plan,
the real target population, accepted precision and loss thresholds and
measured costs. The prospective public/generated comparison can demonstrate
why acquiring those inputs is worthwhile. The exact examples already show
what would be lost by a pooled audit or a premature graph completion; they
do not estimate the institution's benefit from purchasing a prototype.

\section{Use the pooled information when choosing the next audit}
\label{exec17:s7-constraint-design}

The earlier allocation deliberately estimated both stratum moments without
using the exact pooled audit. That is a useful baseline, but the bank should
also ask whether information it already owns changes the next purchase.
Retain precisely the same hypothetical target, the same independent truthful
neighborhood law and the same known response and monetary inputs. Grant the
pooled constraint in \eqref{exec17:s7-pooled-constraint}; no additional outcome
experiment or homogeneous graph restriction is introduced. In this section
the pooled information is already available to every compared method. Its
original acquisition cost is outside the incremental comparison, and would
have to be restored in a complete investment case.

Let $\widehat\theta_z$ be independent stratum audit means. Add any multiple
$\lambda$ of the known zero constraint to the direct value estimator:
\begin{equation}
 \begin{split}
  \widehat v_\lambda
   &=22-\tfrac\lambda4+a_0\widehat\theta_0+a_1\widehat\theta_1,\\
  a_0&=-8+\tfrac34\lambda,\qquad
  a_1=-48+\tfrac14\lambda,\qquad a_0-3a_1=136.
 \end{split}
 \label{exec17:s7-constrained-estimator}
\end{equation}
The estimator is unbiased for every feasible pair of graph moments because
$3\theta_0/4+\theta_1/4=1/4$. This is not an arbitrary change in the target:
all coefficient choices have the same expectation on the supplied information
set. The particular choice $\lambda=0$ recovers the earlier two-mean estimator.
The estimator need not itself lie in the sharp value set in a finite sample;
unbiasedness and feasibility of a realized estimate are distinct properties.

For fixed $\lambda$ and positive audit sizes its variance obeys
\begin{equation}
 \Var(\widehat v_\lambda)
 =\frac{a_0^2\theta_0(1-\theta_0)}{n_0}
  +\frac{a_1^2\theta_1(1-\theta_1)}{n_1}
 \leq \frac14\left(\frac{a_0^2}{n_0}+\frac{a_1^2}{n_1}\right).
 \label{exec17:s7-constrained-variance}
\end{equation}
Use this common one-quarter Bernoulli envelope as a declared design criterion.
It is conservative: the pooled restriction also limits which two actual
variances can occur together. Minimizing the displayed envelope is therefore
a precise and narrower question than finding the best possible estimator or
the actual minimum worst-case variance over the constrained graph family.

Suppose positive prices are $c_0,c_1$ and the incremental audit budget is
$C=c_0n_0+c_1n_1$. For fixed coefficients, Cauchy--Schwarz gives
\begin{equation}
 \frac14\left(\frac{a_0^2}{n_0}+\frac{a_1^2}{n_1}\right)
 \geq \frac{\left(|a_0|\sqrt{c_0}+|a_1|\sqrt{c_1}\right)^2}{4C}.
 \label{exec17:s7-constrained-budget}
\end{equation}
Equality requires $n_z$ proportional to $|a_z|/\sqrt{c_z}$. A zero coefficient
allows zero audits in that stratum; there is then no undefined sample mean in
the estimator. Continuous allocations give the general envelope solution.
Integer allocations need a feasible rounding calculation, except where the
displayed examples attain the bound exactly.

The identity $a_0-3a_1=136$ makes the remaining optimization global rather
than a search over convenient coefficients. If $c_1\leq9c_0$, then
\begin{equation}
 |a_0|\sqrt{c_0}+|a_1|\sqrt{c_1}
 \geq \frac{\sqrt{c_1}}3(|a_0|+3|a_1|)
 \geq \frac{136\sqrt{c_1}}3.
 \label{exec17:s7-stratum1-proof}
\end{equation}
Choose $\lambda=32/3$, so $a_0=0$ and $a_1=-136/3$. Auditing only stratum
one attains both inequalities. Conversely, if $c_1\geq9c_0$,
\begin{equation}
 |a_0|\sqrt{c_0}+|a_1|\sqrt{c_1}
 \geq \sqrt{c_0}(|a_0|+3|a_1|)
 \geq136\sqrt{c_0},
 \label{exec17:s7-stratum0-proof}
\end{equation}
attained by $\lambda=192$, $a_1=0$, $a_0=136$, and auditing only stratum
zero. Thus the joint minimum of the declared envelope is
\begin{equation}
 \mathcal B(C;c_0,c_1)=\frac{4624}{C}
                  \min\left(c_0,\frac{c_1}{9}\right).
 \label{exec17:s7-constrained-optimum}
\end{equation}
At $c_1=9c_0$ the two boundary designs tie. Coefficient choices between the
two kinks also attain the same bound with their corresponding allocations.
The price threshold nine is a conclusion about this envelope and information
set; it is not a universal rule for banking audits.

The stratum-one design has the particularly simple form
\begin{equation}
 \widehat v=\tfrac{58}{3}-\tfrac{136}{3}\widehat\theta_1,
 \qquad
 \widehat\theta_0=\tfrac{1-\widehat\theta_1}{3}.
 \label{exec17:s7-one-stratum}
\end{equation}
The exact pooled information supplies what would otherwise require labels
from the other stratum. This is scientifically different from assuming equal
overlap rates. A realized $\widehat\theta_1>1/2$ gives an infeasible estimated
pair even though the estimator is unbiased. Projecting this sample mean onto
$[0,1/2]$ cannot increase its distance to a true $\theta_1\in[0,1/2]$, and
therefore cannot increase the squared error of the affine value estimate.
That projected estimate need not remain unbiased. Neither operation by itself
creates a confidence interval.

At equal prices and 700 audits, assigning all 700 to stratum one gives the
envelope
\begin{equation}
 \frac{(136/3)^2}{4(700)}=\frac{1156}{1575}\simeq0.7340,
 \label{exec17:s7-pooled-count-comparison}
\end{equation}
in square dollars, compared with $28/25=1.12$ for the earlier optimized
two-mean allocation $(100,600)$. At prices $(1,4)$ and budget 1820, the
integer design $(0,455)$ gives $4624/4095\simeq1.1292$, compared with
$52/35\simeq1.4857$ for $(140,420)$. The target, sampling unit, total count
or budget, and already-available pooled information are matched. These are
exact envelope comparisons, not observed sampling improvements. Finite
enumeration on twelve small Bernoulli audit laws separately checks the
unbiased means and actual variance formula, including the zero-sample boundary
designs.~\citep{exec17_s7_extension}

\subsection{An estimated pooled constraint has a measurable error cost}
\label{exec17:s7-noisy-pool}

An actual pilot is unlikely to deliver its population pooled moment exactly.
Write $\tau=3\theta_0/4+\theta_1/4$, without fixing it at $1/4$.
Eliminating $\theta_0$ from the target gives the identity
\begin{equation}
 v_*=22-\tfrac{32}{3}\tau-\tfrac{136}{3}\theta_1.
 \label{exec17:s7-noisy-value}
\end{equation}
For unbiased estimates of both inputs, the corresponding linear value
estimate has variance
\begin{equation}
 \begin{split}
  \Var(\widehat v_*)&=\tfrac{1024}{9}\Var(\widehat\tau)
       +\tfrac{18496}{9}\Var(\widehat\theta_1)\\
       &\quad+\tfrac{8704}{9}
          \operatorname{Cov}(\widehat\tau,\widehat\theta_1).
 \end{split}
 \label{exec17:s7-noisy-variance}
\end{equation}
Separate independent audit experiments can make the covariance zero; reusing
the same stratum-one labels in the pool does not. A known pool plus one
stratum, two independent estimated summaries, and overlapping summaries are
three different observation laws. The coefficient identity does not excuse
silently replacing one law by another.

If only simultaneous bounds $|\widehat\tau-\tau|\leq d_\tau$ and
$|\widehat\theta_1-\theta_1|\leq d_1$ are warranted, the same event implies
\begin{equation}
 |\widehat v_*-v_*|\leq\tfrac{32}{3}d_\tau+
                            \tfrac{136}{3}d_1.
 \label{exec17:s7-noisy-envelope}
\end{equation}
Their joint probability must be supplied by the sampling argument; two
separately advertised nominal intervals are not automatically simultaneous.
One can intersect the joint region with the probability constraints and map
the intersection into value. Inconsistent exact equalities should instead
trigger an input review. Larger future stratum-one samples do not eliminate
fixed pool bias: ignoring a pool error of magnitude $d$ can retain a monetary
error of $32d/3$ even as those samples grow.

This comparison changes the proposed workflow in a concrete way. Before
pricing the next audit, ask which joint moments are already known, how they
were observed and what covariance or systematic error they carry. If the
pool is defensibly exact, a single strategically selected stratum can suffice
for this target. If it is noisy, preserve its error and reoptimize the joint
design. If it has an unknown truth mechanism, retain the compatible value
set and acquire different information. The source-specific missingness,
linkage and network engines can supply components only after that observation
law is declared; none can turn an uncertain aggregate into an exact scientific
constraint.~\citep{exec17_s7_extension}


\section{Conclusion}
\label{exec17:s7-end}

Observation uncertainty can change a policy's exposure and its transported
value even when separate edge probabilities and ordinary population summaries
are correct. The needed information can be a joint neighborhood moment rather
than more labels of individual fields. A feasible coordinated allocation can
sometimes remove dependence on that moment; a target-specific audit can
sometimes measure it more efficiently. These alternatives require a common
causal request and distinct statements of compatibility, sampling error,
model probability and numerical approximation. The proposed software
composition earns its role by preserving that meaning through the decision.
